% --- ARCHIVO: Capitulos/magia_gema.tex ---

\chapter{La Gema y el Entorno}

La magia en \textit{La Era del Martillo} no es infinita. Es un recurso físico que se extrae de la tierra o se canaliza a través de artefactos antiguos. Comprender el flujo de la energía es la diferencia entre un místico poderoso y un cadáver agotado.

\section{Control de Gema (Gem Control)}
El **Control de Gema** representa la capacidad de un Líder o un portador para canalizar la energía de la Gema Primordial. 

\begin{reglaespecial}{Extracción de Energía}
Cada ronda, el portador de la Gema genera una cantidad de **PE (Puntos de Energía)** igual a su atributo de Control. 
\begin{itemize}
    \item Esta energía puede ser usada por el portador o distribuida a aliados.
    \item Los PE generados por la Gema no se acumulan entre rondas; lo que no se usa, se disipa al final de la activación del último aliado.
\end{itemize}
\end{reglaespecial}

\section{Área de Influencia}
El **Área de Influencia** es el radio táctico (medido en cm desde la peana del Líder) donde su autoridad y energía son tangibles.

\begin{itemize}
    \item \textbf{Canalización:} Las unidades dentro del área pueden usar los PE del Líder como si fueran propios.
    \item \textbf{Transferencia:} El Líder puede gastar una acción para "Cargar" a una unidad aliada. Esa unidad recibe PE que duran toda la ronda, permitiéndole salir del área de influencia y seguir usando habilidades mágicas sin coste adicional de terreno.
\end{itemize}

\section{El Costo del Mundo: Zonas Secas}
Cuando un hechicero o unidad necesita usar una habilidad y no tiene acceso a la energía de una Gema o reserva propia, debe recurrir a la **Extracción Forzada**.

\begin{reglaespecial}{Creación de una Zona Seca}
Si una unidad usa PE del entorno, el punto exacto donde se encuentra se convierte en una \textbf{Zona Seca}.
\begin{enumerate}
    \item Coloca un marcador de 3cm de radio centrado en la unidad actora.
    \item \textbf{Efecto:} Por el resto de la partida, ninguna unidad puede extraer energía del entorno dentro de ese radio.
    \item \textbf{Penalizador:} Las unidades que terminen su movimiento en una Zona Seca sufren un penalizador de -1 a sus tiradas de \textit{Focus}.
\end{enumerate}
\end{reglaespecial}

\section{Habilidades de Gema}
Las habilidades que consumen PE se dividen en:
\begin{description}
    \item[Hechizos Rápidos:] Se activan como Reacción (Coste en PA + PE).
    \item[Rituales:] Requieren una acción completa y suelen tener un alto consumo de PE o requieren sacrificar salud de la unidad.
    \item[Auras:] Habilidades pasivas que se mantienen activas mientras el Líder tenga al menos 1 PE en su reserva.
\end{description}

\vfill
\begin{center}
    \includegraphics[width=0.4\linewidth]{example-image-golden} % Aquí podrías poner un dibujo de la gema
    \\ \textit{"La tierra grita cuando la Gema calla."}
\end{center}